\begin{claim}{D.2.15}{Proved}{Reduced complex projective space generates spectra; the Pr\"ufer ideal at every prime}
Let $D=\Sigma^\infty\mathbb CP^\infty$ be the reduced suspension
spectrum. For every $E_1$-ring spectrum $R$,
\[
 F(D,R)\simeq0\quad\Longrightarrow\quad R\simeq0.
 \tag{D.2.15a}
\]
Consequently
\[
 \begin{gathered}
 \operatorname{loc}(D)=\mathrm{Sp},\\
 \operatorname{loc}\bigl(\Sigma^\infty BC_{p^\infty}\bigr)
   =\operatorname{loc}(D\wedge S/p)=\operatorname{loc}(S/p)
 \end{gathered}
 \tag{D.2.15b}
\]
for every prime $p$, including $2$. In particular the degree-one
Pr\"ufer generator has the full finite $p$-primary thick part.
\end{claim}

\begin{proof}
The first part of D.2.14 proves \textup{(D.2.15a)} whenever the image
of $\eta$ in $\pi_1R$ is zero. We first remove this hypothesis for
derived $2$-complete rings, and then remove the completeness
hypothesis.

\emph{The bottom reduced circle transfer is an odd multiple of $\nu$.}
Use the Euler--Thom sequence \textup{(D.2.14c)} and split off the
basepoint of $\mathbb CP^\infty_+$:
\[
 S^{-2}\longrightarrow\operatorname{Th}(-L)
  \longrightarrow S\vee D\xrightarrow{\delta}S^{-1}.
\]
Let $t:D\to S^{-1}$ be the restriction of $\delta$, let
$i:S^2\to D$ include the bottom cell, and put
\[
 b=\Sigma(ti):S^3\longrightarrow S.
 \tag{D.2.15c}
\]
We will show that $b=k\nu$ in
$\pi_3S_{(2)}\cong\mathbb Z/8$ with $k$ odd. Here $\nu$ is the
stable quaternionic Hopf class.

There are natural maps of homotopy fibers
\[
 F_{ti}\longrightarrow F_t\longrightarrow F_\delta
       \simeq\operatorname{Th}(-L),
 \qquad F_{ti}\simeq\Sigma^{-2}Cb,
 \tag{D.2.15d}
\]
where $Cb=\operatorname{cofib}(b:S^3\to S)$.
Each of the two maps in \textup{(D.2.15d)} induces an isomorphism
on mod-$2$ cohomology in degrees $-2$ and $2$. This follows directly
from the fiber-sequence cohomology exact sequences: the maps
$S^2\to D\to S\vee D$ are isomorphisms in degree $2$, all their
sources and targets have zero cohomology in degrees $-2,-1$, and
the common target $S^{-1}$ supplies the degree-$-2$ fiber class.
The extra basepoint summand changes only degree $0$ here.

Let $u\in H^{-2}(\operatorname{Th}(-L);\mathbb F_2)$ be the
virtual Thom class and let $x=c_1(L)\bmod2$. The Thom identity and
the Whitney product formula give
\[
 w(-L)=(1+x)^{-1}=1+x+x^2+\cdots,
 \qquad \operatorname{Sq}^4u=x^2u\ne0.
 \tag{D.2.15e}
\]
The ordinary Thom identity is given in
\cite[\S8, pp.~91--93]{EH5MilnorStasheff1974}.
Its use for $-L$ requires no additional orientation assertion:
over $\mathbb CP^2$, choose a finite-dimensional complement $L^\perp$
to $L$ in a trivial complex bundle, apply the ordinary identity to
$L^\perp$, and desuspend by the dimension of that trivial bundle.
The Whitney product formula gives exactly the inverse total class
in \textup{(D.2.15e)}.

Pulling \textup{(D.2.15e)} through \textup{(D.2.15d)} shows that
$\operatorname{Sq}^4$ is nonzero from the bottom to the top class
of $Cb$. Thus $b$ has mod-$2$ Hopf invariant one. Since
$\pi_3S_{(2)}=\mathbb Z/8\{\nu\}$, this says precisely that its
coefficient $k$ is odd; see
\cite[Example 7.17, pp.~45--46, and equation (8.1)]{EH5Adams1966}.
Only this parity is needed, so no choice of sign or normalization
of the circle transfer is involved.

\emph{A derived $2$-complete ring.}
Suppose now that $R$ is derived $2$-complete and $F(D,R)=0$.
The composite of $t$ with the shifted unit
$S^{-1}\to\Sigma^{-1}R$ is null. Its bottom restriction makes
$b_R=0$. All odd integers act invertibly on a derived $2$-complete
spectrum, so \textup{(D.2.15c)--(D.2.15e)} imply
\[
 \nu_R=0.
 \tag{D.2.15f}
\]

Let $C\eta=\operatorname{cofib}(\eta:S^1\to S)$ and set
$M=R\wedge C\eta$, viewed as a left $R$-module. The stable action
of $\eta$ on $C\eta$ factors as
\[
 \eta_{C\eta}:\Sigma C\eta\longrightarrow S^3
       \xrightarrow{\nu}S\longrightarrow C\eta;
 \tag{D.2.15g}
\]
this is \cite[equation (8.4), p.~1279]{EH5AngeltveitRognes2005}.
For the implication needed here one can also derive a possibly
non-normalized version directly. The composite with the suspended
bottom inclusion is zero, so $\eta_{C\eta}$ factors through $S^3$.
The map $\pi_3S\to\pi_3C\eta$ is surjective, since multiplication
by $\eta$ from $\pi_1S$ to $\pi_2S$ is an isomorphism. Thus, at
$2$, the remaining factor comes from a multiple of $\nu$.
Either description and \textup{(D.2.15f)} give
$\eta_M=0$ as an $R$-linear map.

Consider the unital $E_1$-ring
\[
 A=\operatorname{End}_R(M).
\]
Because $C\eta$ is finite, its underlying spectrum is
\[
 A\simeq(C\eta)^\vee\wedge R\wedge C\eta.
 \tag{D.2.15h}
\]
The class of spectra $Z$ with $F(D,Z)=0$ is stable and closed
under finite colimits and retracts. It therefore contains
$R\wedge K$ for every finite spectrum $K$. In particular
\textup{(D.2.15h)} gives $F(D,A)=0$.
The image of $\eta$ under the unit of $A$ is exactly the
$R$-linear self-map $\eta_M$, so $\eta_A=0$.
D.2.14(1), which applies to $E_1$-rings, now gives $A=0$.
Its unit represents $\operatorname{id}_M$, hence $M=0$.

It follows that $\eta:\Sigma R\to R$ is an equivalence, since its
cofiber is $M$. But $\eta^4=0$ in the stable sphere; this relation
is recalled explicitly in
\cite[proof of Lemma 8.11, p.~1285]{EH5AngeltveitRognes2005}.
An invertible nilpotent action forces $R=0$. This proves
\textup{(D.2.15a)} for derived $2$-complete $E_1$-rings.

\emph{An arbitrary ring.}
Let $R$ now be any $E_1$-ring with $F(D,R)=0$, and put
$C=\operatorname{fib}(S\to S[1/2])$.
As established in the local-unit argument of D.2.14,
\[
 \widehat R_2\simeq F(C,R),\qquad
 \widehat R_2/2\simeq R/2,
 \tag{D.2.15i}
\]
and derived $2$-completion is a monoidal localization, so
$\widehat R_2$ inherits an $E_1$-ring structure. Also
\[
 F(D,\widehat R_2)\simeq F(D\wedge C,R)=0.
\]
Indeed $\operatorname{loc}(D)$ is a tensor ideal, and $R$ is right
orthogonal to that whole localizing subcategory. The completed
case makes $\widehat R_2=0$. By \textup{(D.2.15i)}, $R/2=0$,
so multiplication by $2$ is invertible on $R$. Since $2\eta=0$,
we have $\eta_R=0$; D.2.14(1) gives $R=0$.

\emph{From ring detection to actual localizing generation.}
The subcategory $\mathcal I=\operatorname{loc}(D)$ is generated by
one object and is a stable tensor ideal. Its accessible exact
localization therefore has a localized symmetric monoidal
structure; the inclusion of local objects is lax symmetric
monoidal. In particular the underlying spectrum of the localized
unit, $L_{\mathcal I}S$, is a unital $E_\infty$-ring and is
$D$-null; see
\cite[Example 2.2.1.7, Proposition 2.2.1.9]{FGHLurieHA}.
The detector \textup{(D.2.15a)} gives $L_{\mathcal I}S=0$.
Thus $S\in\mathcal I$, and $\mathcal I=\mathrm{Sp}$.
This argument does not require the localization to be smashing.

Finally smash a localizing construction of $S$ from $D$ with
$S/p$. The functor $-\wedge S/p$ is exact and preserves all
colimits, so it proves
$S/p\in\operatorname{loc}(D\wedge S/p)$.
The reverse inclusion follows from the tensor-ideal property of
$\operatorname{loc}(S/p)$. The Pr\"ufer comparison
\[
 \operatorname{loc}(\Sigma^\infty BC_{p^\infty})
       =\operatorname{loc}(D\wedge S/p)
\]
was proved for every prime in D.2.14. This completes
\textup{(D.2.15b)}.
\end{proof}

For historical comparison, the equality of smash Bousfield classes
$\langle\mathbb CP^\infty\rangle=\langle S\rangle$ was already
reported as a theorem of Hopkins in the discussion preceding
\cite[Conjecture 8.7]{JHoveyPalmieri1999}.
The proof above establishes localizing generation directly; it
does not assume the general comparison between Bousfield order
and localizing membership discussed in that paper's Section~9.
